\documentclass[
  12pt,
  a4paper,
  oneside,
  parskip=half,
  headings=normal
]{scrreprt}

% -----------------------------------------------------------------------------
% Sprache, Typografie und Seitenlayout
% -----------------------------------------------------------------------------
\usepackage[T1]{fontenc}
\usepackage[utf8]{inputenc}
\usepackage[ngerman]{babel}
\usepackage{lmodern}
\usepackage{microtype}
\usepackage{geometry}
\usepackage{setspace}
\usepackage{enumitem}
\usepackage{xcolor}
\usepackage{listings}
\usepackage{graphicx}
\usepackage[hidelinks,hypertexnames=false]{hyperref}

\geometry{
  left=3cm,
  right=2.5cm,
  top=2.5cm,
  bottom=2.5cm
}
\onehalfspacing
\setlist{itemsep=0.25em, topsep=0.5em}

% -----------------------------------------------------------------------------
% Metadaten: vor der Abgabe ausfüllen
% -----------------------------------------------------------------------------
\newcommand{\titel}{SRD Arena}
\newcommand{\untertitel}{Entwicklung einer Engine für taktische Kämpfe nach dem SRD~5.2.1}
\newcommand{\autor}{Vorname Nachname}
\newcommand{\matrikelnummer}{00000000}
\newcommand{\studiengang}{Studiengang}
\newcommand{\modul}{Modulname}
\newcommand{\dozent}{Name der Lehrperson}
\newcommand{\hochschule}{Name der Hochschule}
\newcommand{\abgabedatum}{\today}

% -----------------------------------------------------------------------------
% Hilfsmittel für Entwurfsnotizen und Quelltext
% -----------------------------------------------------------------------------
\definecolor{hinweisfarbe}{RGB}{120,70,0}
\definecolor{codehintergrund}{RGB}{247,247,247}

\newcommand{\vorgabe}[1]{%
  \begin{quote}
    \small\color{hinweisfarbe}\textbf{Vorgabe:} #1
  \end{quote}
}

\newcommand{\todo}[1]{%
  \par\noindent\textcolor{red!70!black}{\textbf{TODO:} #1}\par
}

\lstdefinestyle{projektcode}{
  basicstyle=\ttfamily\small,
  backgroundcolor=\color{codehintergrund},
  frame=single,
  breaklines=true,
  columns=fullflexible,
  keepspaces=true,
  showstringspaces=false
}
\lstset{style=projektcode}

\begin{document}

% -----------------------------------------------------------------------------
% Titelseite
% -----------------------------------------------------------------------------
\begin{titlepage}
  \centering
  {\Large \hochschule\par}
  \vspace{1.5cm}
  {\large \studiengang\par}
  {\large \modul\par}
  \vfill
  {\Huge\bfseries \titel\par}
  \vspace{0.5cm}
  {\Large \untertitel\par}
  \vfill
  \begin{tabular}{ll}
    Autor: & \autor \\
    Matrikelnummer: & \matrikelnummer \\
    Lehrperson: & \dozent \\
    Abgabedatum: & \abgabedatum
  \end{tabular}
  \vfill
\end{titlepage}

\pagenumbering{Roman}
\tableofcontents
\clearpage
\pagenumbering{arabic}

% -----------------------------------------------------------------------------
\chapter{Einleitung}
% -----------------------------------------------------------------------------

Dungeons \& Dragons (D\&D) ist ein beliebtes Pen-and-Paper-Rollenspiel, in
dem taktische, rundenbasierte Kämpfe (Encounter) einen großen Teil des Spiels
ausmachen. Üblicherweise kämpft ein Team von Spielern gegen ein Team von NPCs
(Non-Player Characters), die vom Spielleiter gesteuert werden.

Das System Reference Document (SRD) 5.2.1 ist eine unter der
Creative-Commons-Lizenz CC BY 4.0 verfügbare Fassung der Spielregeln. Sie
enthält den größten Teil der Grundmechaniken des Spiels, jedoch nicht sämtliche
Inhalte der proprietären Regelwerke. So kann ein Spieler beispielsweise eine
Klasse und Subklasse wählen; im SRD steht aber pro Klasse nur eine Subklasse
zur Verfügung, während das proprietär lizenzierte Spielerhandbuch weitere
Subklassen enthält.

SRD Arena setzt zunächst einen Teil des Kampfsystems des SRD~5.2.1 als
rundenbasiertes 2D-Spiel um. Mehrere vordefinierte Kampfszenarien demonstrieren
die implementierten Funktionen. Mithilfe von Vorlagen können Nutzer weitere
Szenarien in strukturierten Dateien definieren und anschließend spielen.

\section{Motivation}

Die meisten Spielleiter entwerfen Kampfszenarien bereits vor dem eigentlichen
Spieltermin. Ihr Ziel ist es, den Spielern neben einem interessanten Narrativ
auch eine taktische Herausforderung zu bieten, ohne den Kampf unangemessen
schwer zu gestalten.

Im Rahmen der umgesetzten Regeln können Spielleiter mit SRD Arena ihre
Kampfszenarien vor dem Spiel ausprobieren. Viele Berechnungen und Würfelwürfe,
die sonst manuell erfolgen müssten, werden dabei automatisiert. Mittelfristig,
jedoch nach Abgabe des Projekts, soll der gesamte Umfang der SRD-Regeln
abgedeckt werden.

Langfristig könnte das Projekt um Bots erweitert werden, die die
Kampfentscheidungen aller Teilnehmer automatisieren. Kämpfe könnten dann
vielfach simuliert und die entstehenden Daten in robuste Kennzahlen zur
Ausgewogenheit eines Szenarios überführt werden. Diese Funktionalität stellt
die weiterführende Motivation des Projekts dar, kann aufgrund der begrenzten
Bearbeitungszeit jedoch erst nach der Abgabe umgesetzt werden.

\section{Zielsetzung}

SRD Arena ist eine Engine für taktische Kämpfe auf einem 2D-Spielfeld nach dem
SRD-5.2.1-Regelwerk. Charaktere werden wahlweise durch Nutzereingaben oder
durch einfache Bots gesteuert. Szenarien werden in strukturierten Textdateien
definiert; dafür stehen eine Vorlage und mehrere Demonstrationsszenarien zur
Verfügung.

\subsection{Verwendete Pakete}

\paragraph{Kernabhängigkeiten}
\begin{itemize}
  \item \texttt{pydantic}: Validierung der JSON-Daten
  \item \texttt{PySide6}: grafische Benutzeroberfläche
  \item \texttt{scipy}: geometrische Berechnungen
\end{itemize}

\paragraph{Entwicklungsabhängigkeiten}
\begin{itemize}
  \item \texttt{hypothesis}: Property-Based Testing
  \item \texttt{mypy}: statische Typprüfung
  \item \texttt{pytest}: Unit-Tests
  \item \texttt{pytest-cov}: Messung der Testabdeckung
  \item \texttt{ruff}: Linter und Formatter
\end{itemize}

\subsection{Abgrenzung des MVP}

Wesentlich ist eine stabile Encounter-Engine für vorhandene Spieldaten. Der
MVP umfasst ein klar abgegrenztes Subset der SRD-Kampfmechaniken.

\paragraph{Kampfablauf}
\begin{itemize}
  \item Initiative
  \item Zugstruktur
  \item Action, Bonus Action und Movement
  \item Opportunity Attack als erste Reaktionsmechanik
\end{itemize}

\paragraph{Bewegung}
\begin{itemize}
  \item rasterbasierte Bewegung
\end{itemize}

\paragraph{Kampfmechaniken}
\begin{itemize}
  \item Nahkampf- und Fernkampfangriffe
  \item Zauber, deren zugrunde liegende Mechaniken unterstützt werden
  \item Attack Rolls und Saving Throws
  \item Armor Class und Critical Hits
  \item Waffen mit unterschiedlichen Schadenswürfeln
\end{itemize}

\paragraph{Effekte}
\begin{itemize}
  \item Radius- und Kegel-Flächeneffekte
  \item Heilung
  \item einfache Zustände, insbesondere \emph{Blinded} und \emph{Prone}
\end{itemize}

\paragraph{Nicht Teil des MVP}
\begin{itemize}
  \item vollständige SRD-Abdeckung
  \item vollständige Zauberliste
  \item vollständige Klassenprogression
  \item vollständige Monster-KI
\end{itemize}

\paragraph{Erfolgskriterien}
\begin{itemize}
  \item Szenarien und Spieldaten können aus JSON geladen werden.
  \item Ungültige Eingaben werden durch Schema-Validierung erkannt.
  \item Ein vollständiger Encounter lässt sich ausführen; anschließend wird
        das Textabenteuer fortgesetzt.
  \item Derselbe Seed erzeugt denselben Encounter-Verlauf.
  \item Ein Baseline-Agent spielt einen vollständigen Encounter mit
        ausschließlich legalen Aktionen.
  \item Aktionen, Würfe, Zustandsänderungen und Kampfende sind vollständig im
        Log enthalten.
  \item Alle unterstützten Mechaniken sind durch automatisierte Tests
        abgedeckt.
\end{itemize}

\paragraph{Optionale Erweiterungen}
\begin{itemize}
  \item weitere AoE-Typen und Flächeneffekte über mehrere Runden
  \item komplexere Zustände
  \item Summons, die neue Akteure erzeugen und in die Initiative aufnehmen
  \item weitere Reaktionsmechaniken
  \item Konzentration, Grapple und Shove
  \item Vision, Licht, Hindernisse und Line of Sight
  \item bessere Gegner-Agenten
  \item aufwendigere GUI-Animationen
\end{itemize}

\subsection{Nächste Entscheidungen zur Implementierungsstrategie}
\begin{itemize}
  \item Feature-Abhängigkeitsgraph und Abhängigkeiten zwischen Mechaniken
        dokumentieren
  \item spätere Erweiterungen priorisieren
  \item JSON-Schemas finalisieren
  \item Formate für Spieldaten, Szenarien und Traces festlegen
  \item Validierungsregeln definieren
  \item Schnittstelle zwischen Kernlogik und GUI festlegen
  \item Regelentscheidungen und Zustandsverwaltung in der Kernlogik belassen
  \item GUI auf Darstellung und Eingabemöglichkeiten begrenzen
\end{itemize}

\subsection{Meilensteine}
\begin{itemize}
  \item Kernlogik aufräumen
  \item JSON-Schemas finalisieren
  \item Schnittstelle für Aktionen und Zielauswahl fertigstellen
  \item Determinismus sicherstellen
  \item Testabdeckung für alle MVP-Mechaniken herstellen
  \item Baseline-Agent stabilisieren
  \item GUI anbinden
  \item optionale Erweiterungen umsetzen
\end{itemize}

\subsection{Erwartete Schwierigkeiten}
\begin{itemize}
  \item Reaktionen verkomplizieren die Zugstruktur, da andere Teilnehmer
        während eines fremden Zuges kurzfristig handeln können.
  \item AoE-Effekte erzeugen zahlreiche Randfälle; im MVP werden nur
        ausgewählte AoE-Typen unterstützt.
  \item JSON-Daten können formal gültig sein, obwohl die benötigte Mechanik
        noch nicht implementiert ist. Schema-Validierung und Laufzeitfehler
        müssen deshalb getrennt behandelt werden.
  \item Viele Mechaniken bauen aufeinander auf. Fehlende Grundmechaniken können
        mehrere spätere Erweiterungen blockieren.
  \item GUI und Gegner-Agenten müssen mit wachsendem Funktionsumfang angepasst
        werden.
\end{itemize}

% -----------------------------------------------------------------------------
\chapter{Theoretische Grundlagen und Werkzeuge}
% -----------------------------------------------------------------------------

\section{Technologien}
\todo{Die eingesetzten Technologien einordnen und ihre Auswahl begründen.}

\section{Bibliotheken und Pakete}

\vorgabe{Beschreiben, welche Pakete und Bibliotheken zum Einsatz kamen.}

\begin{itemize}
  \item Typprüfung: \texttt{mypy}
  \item Linting und Formatierung: \texttt{ruff}
  \item Tests: \texttt{pytest}
  \item Property-Tests: \texttt{hypothesis}
  \item Testabdeckung: \texttt{pytest-cov}
  \item Datenvalidierung: \texttt{pydantic}
  \item Paketverwaltung: \texttt{uv}
  \item Konfiguration: \texttt{pyproject.toml}
  \item Debugging und Replay: Logging und JSON-Traces
  \item Nicht-Python-Datenquelle: JSON-Daten der SRD-Regeln von 5e.tools als
        Grundlage für den authored content
\end{itemize}

\section{Mensch und KI}

\vorgabe{Die Rolle der KI-Unterstützung im Projekt erläutern.}

In diesem Projekt wurde ein großer Teil der Programmlogik mit
KI-Unterstützung generiert. Viele Strukturideen wurden mithilfe der KI auf
Lücken geprüft und iterativ angepasst. Außerdem wurden Regelabschnitte, die in
JSON-Dateien zunächst nur als Prosa vorlagen, mithilfe der KI in strukturierte
JSON-Definitionen übertragen.

% -----------------------------------------------------------------------------
\chapter{Konzept und Implementierung}
% -----------------------------------------------------------------------------

\section{Programmstruktur}
\todo{Die obersten Pakete und ihre Verantwortlichkeiten anhand einer Grafik
oder eines kompakten Paketbaums erläutern.}

\section{Kernlogik}
\todo{Zustandsmodell, Zugablauf und Ausführung von Aktionen beschreiben.}

\section{Datenverarbeitung}
\todo{Laden, Validieren, Indizieren und Bauen der Inhaltsdaten beschreiben.}

\section{Architektur- und Designentscheidungen}

\vorgabe{Begründen, warum die Software auf diese Weise strukturiert wurde.
Dazu gehören beispielsweise die Wahl bestimmter Entwurfsmuster,
objektorientierter Strukturen oder funktionaler Ansätze sowie verworfene
Alternativen und deren Gründe.}

Wesentliche Entscheidungen sind:
\begin{itemize}
  \item authored content in JSON-Dateien
  \item Trennung in Content-, Domain-, Runtime- und GUI-Schicht
\end{itemize}

\subsection{Content-Schicht}

\begin{lstlisting}[caption={Vom Dateisystem zum Domain-Zauber}]
Dateisystem
    | load_spell_catalog()
    v
SpellCatalog mit validierten SpellSchema-Objekten
    | find()
    v
SpellSchema
    | build_spell()
    v
Domain-Objekt Spell
\end{lstlisting}

\begin{lstlisting}[caption={Validierung und Indizierung}]
JSON-Datei
    | parsen und validieren
    v
SpellSchema
    | speichern und indizieren
    v
SpellCatalog
\end{lstlisting}

\section{KI-Konsultation bei Designentscheidungen}

\vorgabe{Darstellen, ob die KI bereits bei Planung und Architektur konsultiert
wurde, welche Vorschläge übernommen wurden und wo bewusst von Empfehlungen
abgewichen wurde.}

Die KI unterstützte insbesondere beim Erlernen und Anwenden einer geschichteten
Architektur, beim Schemaentwurf für Inhaltsdaten und bei der Recherche zu
Regelgrenzfällen wie \emph{Calm Emotions} und Immunitäten. Die Vorschläge
wurden nicht unverändert übernommen:

\begin{enumerate}
  \item Bei einer frühen Zauberimplementierung wurde gemeinsam mit der KI ein
        Fahrplan erstellt. Der Code wurde jedoch schnell unübersichtlich. Eine
        Empfehlung, zunächst weitere repräsentative Zauber zu implementieren,
        wurde verworfen, um die problematische Struktur nicht weiter
        auszubauen.
  \item In der experimentellen Phase existierten nur ein Spieler und ein
        Gegner. Daraus entstand die Annahme einer primären Kreatur und parallel
        entwickelte Logik für Spieler, Verbündete und Gegner. Diese historisch
        gewachsene Struktur musste später vereinheitlicht werden.
  \item Bei Zuständen verhinderte Domänenwissen eine falsche Abstraktion:
        \emph{Unconscious} impliziert \emph{Incapacitated} und verursacht
        \emph{Prone}. Endet \emph{Unconscious}, endet das davon abhängige
        \emph{Incapacitated}; \emph{Prone} kann dagegen unabhängig bestehen
        bleiben.
  \item Regelprosa wird nicht zur Laufzeit geparst. Stattdessen werden
        ausführbare Mechaniken in JSON angereichert. Dadurch bleibt weniger
        Parsing-Logik im Programmcode verborgen.
  \item Bei generischen Kreaturentyp-Anforderungen und der Modellierung von
        Immunität gegenüber der Unterdrückung eines Zustands wurden erste
        KI-Vorschläge anhand des Regeltexts überprüft und korrigiert.
\end{enumerate}

\section{Code-Generierung}

\vorgabe{Erläutern, welche Bestandteile wie programmiert wurden.}

Für die Entwicklung wurde OpenAI Codex eingesetzt.

\subsection{KI-generierte Bestandteile}

\vorgabe{Aufführen, welche Module, Funktionen oder Tests maßgeblich mit
KI-Unterstützung erstellt wurden, und begründen, weshalb sich der KI-Einsatz
dafür eignete.}

\begin{itemize}
  \item Übertragung verbleibender Regelprosa in strukturierte JSON-Daten für
        Zauber
  \item schrittweise Implementierung von Regeln wie Trefferpunkten, Angriffen
        und Zaubern
\end{itemize}

\todo{Konkrete Module, Funktionen und Tests mit Commit- oder Dateiverweisen
ergänzen.}

\subsection{Selbst programmierte Bestandteile}

\vorgabe{Kernlogiken, komplexe Workarounds oder projektspezifische Funktionen
beschreiben, die vollständig selbst programmiert wurden, und die Notwendigkeit
der eigenen Programmierleistung begründen.}

\begin{itemize}
  \item eigenständige Bewertung und Korrektur der Repository-Struktur
  \item fachliche Prüfung der von der KI vorgeschlagenen Regelmodelle
\end{itemize}

\todo{Konkrete selbst programmierte Bestandteile ergänzen.}

\section{Prompt-Engineering und Iteration}

\vorgabe{Die Zusammenarbeit mit der KI beschreiben. Insbesondere erläutern, ob
generierter Code direkt übernommen oder in kritischen Iterationsschleifen
geprüft und überarbeitet wurde. Dazu gehören manuelle Korrekturen von
Halluzinationen und veralteter Syntax.}

\todo{Ein bis zwei konkrete Prompt- und Überarbeitungszyklen dokumentieren.}

% -----------------------------------------------------------------------------
\chapter{Ergebnisse}
% -----------------------------------------------------------------------------

\section{Funktionsnachweis}
\todo{Testergebnisse, unterstützte Mechaniken und reproduzierbare
Qualitätsprüfungen dokumentieren.}

\section{Beispieldurchlauf}
\todo{Ein Szenario auswählen und den Ablauf anhand von Screenshots oder
Logauszügen darstellen.}

% -----------------------------------------------------------------------------
\chapter{Diskussion und Fazit}
% -----------------------------------------------------------------------------

\section{Zielerreichung}
\vorgabe{Bewerten, inwiefern das gesetzte Ziel erreicht wurde.}
\todo{Erfolgskriterien aus Abschnitt 1.2 einzeln bewerten.}

\section{Herausforderungen und kritische Reflexion}
\todo{Technische und organisatorische Herausforderungen zusammenfassen.}

\section{Abweichungen vom ursprünglichen Implementierungsplan}
\todo{Abweichungen, Ursachen und Auswirkungen dokumentieren.}

\section{Kritische Würdigung}

\vorgabe{Reflektieren, was gut lief, welche unerwarteten technischen Probleme
oder Schwierigkeiten bei der Zusammenarbeit mit der KI auftraten und was bei
einem zukünftigen Projekt anders gemacht würde.}

\todo{Stärken, Schwächen und konkrete Verbesserungen gegenüberstellen.}

\section{Ausblick}
\vorgabe{Der Ausblick darf auch eine subjektive Einschätzung enthalten.}
\todo{Nächste technische Schritte und mögliche Weiterentwicklung erläutern.}

% -----------------------------------------------------------------------------
\chapter{Literatur- und Quellenverzeichnis}
% -----------------------------------------------------------------------------

\vorgabe{Dieser Abschnitt ist nur erforderlich, wenn zitierfähige Quellen
verwendet wurden.}

\begin{itemize}
  \item \todo{System Reference Document 5.2.1 bibliografisch vollständig
        angeben.}
  \item \todo{Dokumentationen und weitere zitierte Quellen ergänzen.}
\end{itemize}

% -----------------------------------------------------------------------------
\chapter{Anhang}
% -----------------------------------------------------------------------------

\vorgabe{In den Anhang gehören Inhalte, die für den Haupttext zu umfangreich
sind.}

\section{Bedienungsanleitung}
\todo{Installation, Programmstart, Szenarioauswahl und Bedienung dokumentieren.}

\section{Ausgewählte Code-Snippets}
\vorgabe{Die ausgewählten Code-Snippets sollten jeweils nicht länger als eine
Seite sein.}

% Beispiel:
% \lstinputlisting[
%   language=Python,
%   firstline=1,
%   lastline=30,
%   caption={Beispiel für eine zentrale Implementierung}
% ]{src/srd_arena/example.py}

\todo{Repräsentative Codeausschnitte auswählen und erläutern.}

\end{document}
